\documentclass[12pt,a4paper]{article}

% Краткое сообщение Ярославу:
% В заданиях 7 и 8 базовые промежутки выбраны правильно, но в окончательной записи потеряна периодичность.
% Задание 6 по присланному фото не удалось надёжно прочитать целиком, поэтому оно не включено в ревью.

\usepackage[a4paper,margin=2cm]{geometry}
\usepackage[T2A]{fontenc}
\usepackage[utf8]{inputenc}
\usepackage[russian]{babel}
\usepackage{amsmath,amssymb,array,booktabs,longtable,xcolor,hyperref,tikz,fancyhdr}
\hypersetup{colorlinks=true,linkcolor=blue!55!black,urlcolor=blue!55!black}
\setlength{\parindent}{0pt}
\setlength{\parskip}{0.55em}
\renewcommand{\arraystretch}{1.2}
\pagestyle{fancy}\fancyhf{}\fancyfoot[C]{\small Лёвин Артём Александрович эксклюзивно для Ярослава Гаврилова}\renewcommand{\headrulewidth}{0pt}\renewcommand{\footrulewidth}{0pt}
\newcommand{\err}{\textcolor{red}{Ошибка:}}
\newcommand{\idea}{\textcolor{blue!60!black}{Ключевая идея:}}
\newcommand{\result}{\textcolor{teal!60!black}{Итог:}}
\begin{document}
\thispagestyle{empty}
\vspace*{\fill}
\begin{center}{\LARGE\bfseries Ревью домашней работы\par}\vspace{1.1cm}{\large Автор: Лёвин Артём Александрович\par}\vspace{0.6cm}{\large 3 октября 2026 года\par}\end{center}
\vspace*{\fill}
\begin{center}\begin{tikzpicture}\draw[blue!35!black,line width=.7pt] (0,0)..controls(2,.7)and(4,-.7)..(6,0);\end{tikzpicture}\end{center}
\newpage
{\small\setlength{\tabcolsep}{2pt}
\begin{longtable}{>{\raggedright\arraybackslash}p{0.12\linewidth}>{\raggedright\arraybackslash}p{0.20\linewidth}>{\raggedright\arraybackslash}p{0.30\linewidth}>{\raggedright\arraybackslash}p{0.15\linewidth}>{\raggedright\arraybackslash}p{0.15\linewidth}}
\toprule № задания & Тема & Суть ошибки & Ваш ответ & Правильный ответ\\\midrule\endfirsthead
\toprule № задания & Тема & Суть ошибки & Ваш ответ & Правильный ответ\\\midrule\endhead
\hyperref[task:7]{7} & Тригонометри\-ческое неравенство & Найден только один базовый промежуток; периодические сдвиги на целое число полных оборотов не записаны. & $[-\pi/4;\pi/4]$ & $[-\pi/4+2\pi n;$\newline$\pi/4+2\pi n]$\\\midrule
\hyperref[task:8]{8} & Тригонометри\-ческое неравенство & Найден только один базовый промежуток; отсутствует общая периодическая запись всех решений. & $[\pi/3;5\pi/3]$ & $[\pi/3+2\pi n;$\newline$5\pi/3+2\pi n]$\\\bottomrule
\end{longtable}}
\newpage
\thispagestyle{empty}\vspace*{\fill}\begin{center}{\LARGE\bfseries Разбор ошибок}\end{center}\vspace*{\fill}\newpage

\phantomsection\label{task:7}
\section*{Задание №7}
\textbf{Текст задания}

Решить неравенство
\[
\cos x\ge \frac{\sqrt2}{2}.
\]

\textbf{Ваше решение}

По тригонометрической окружности отмечена дуга, на которой косинус не меньше половины корня из двух, и записан промежуток
\[
x\in\left[-\frac\pi4;\frac\pi4\right].
\]

\textbf{Ваш ответ}
\[
x\in\left[-\frac\pi4;\frac\pi4\right].
\]

\textbf{\err\ ответ записан только для одного оборота окружности.}

\textbf{Где именно возникает ошибка}

Последний правильный шаг --- определение базовой дуги от минус одной четверти числа пи до одной четверти числа пи. Первый неполный шаг --- запись этой дуги как окончательного ответа без учёта периодичности косинуса.

\textbf{Почему это неверно}

Косинус имеет период два числа пи. Если некоторое значение аргумента удовлетворяет неравенству, то после прибавления или вычитания любого целого числа полных периодов косинус не изменится. Поэтому решение не может ограничиваться только одним промежутком около нуля.

\textbf{\idea} сначала на одном периоде находим нужную дугу, а затем ко всем её точкам добавляем целое число периодов, то есть два числа пи, умноженные на целое число.

\textbf{Исправление}

Базовая дуга найдена верно:
\[
-\frac\pi4\le x\le \frac\pi4.
\]
Теперь добавляем период два числа пи:
\[
-\frac\pi4+2\pi n\le x\le \frac\pi4+2\pi n,
\qquad n\in\mathbb Z.
\]

\textbf{Полное правильное решение}

Равенство
\[
\cos x=\frac{\sqrt2}{2}
\]
на одном полном обороте достигается при углах одна четверть числа пи и минус одна четверть числа пи. Косинус не меньше указанного значения на дуге между этими точками, содержащей ноль. Поэтому на одном периоде получаем
\[
x\in\left[-\frac\pi4;\frac\pi4\right].
\]
Так как период косинуса равен двум числам пи, все решения имеют вид
\[
x\in\left[-\frac\pi4+2\pi n;\frac\pi4+2\pi n\right],
\qquad n\in\mathbb Z.
\]

\textbf{Проверка}

При $x=0$ косинус равен единице, поэтому точка действительно входит в решение. При $x=\pi/2$ косинус равен нулю, поэтому эта точка в решение не входит. После прибавления двух чисел пи значения косинуса повторяются, что подтверждает периодическую запись.

\textbf{\result}
\[
\boxed{x\in\left[-\frac\pi4+2\pi n;\frac\pi4+2\pi n\right],\quad n\in\mathbb Z.}
\]

\textbf{Задача на закрепление}

Решите неравенство
\[
\cos x\ge \frac12.
\]

\newpage
\phantomsection\label{task:8}
\section*{Задание №8}
\textbf{Текст задания}

Решить неравенство
\[
\cos x\le \frac12.
\]

\textbf{Ваше решение}

По тригонометрической окружности отмечена дуга, на которой косинус не превосходит одной второй, и записан базовый промежуток
\[
x\in\left[\frac\pi3;\frac{5\pi}3\right].
\]

\textbf{Ваш ответ}
\[
x\in\left[\frac\pi3;\frac{5\pi}3\right].
\]

\textbf{\err\ отсутствует общая периодическая запись решения.}

\textbf{Где именно возникает ошибка}

Последний правильный шаг --- выбор на окружности дуги от одной трети числа пи до пяти третей числа пи. Первый неполный шаг --- завершение решения этим единственным промежутком.

\textbf{Почему это неверно}

Неравенство должно быть решено для всех действительных значений аргумента. Косинус повторяет свои значения через два числа пи, поэтому тот же допустимый промежуток повторяется после каждого целого числа полных оборотов как вправо, так и влево.

\textbf{\idea} базовый промежуток нужно переносить на каждый период функции, прибавляя два числа пи, умноженные на целое число.

\textbf{Исправление}

К найденному промежутку добавляем период:
\[
\frac\pi3+2\pi n\le x\le \frac{5\pi}3+2\pi n,
\qquad n\in\mathbb Z.
\]

\textbf{Полное правильное решение}

Равенство
\[
\cos x=\frac12
\]
на промежутке от нуля до двух чисел пи достигается при углах одна треть числа пи и пять третей числа пи. Между этими точками, проходя через число пи, косинус не превосходит одной второй. Значит, на одном периоде
\[
x\in\left[\frac\pi3;\frac{5\pi}3\right].
\]
С учётом периода два числа пи получаем
\[
x\in\left[\frac\pi3+2\pi n;\frac{5\pi}3+2\pi n\right],
\qquad n\in\mathbb Z.
\]

\textbf{Проверка}

При $x=\pi$ косинус равен минус единице, поэтому точка входит в решение. При $x=0$ косинус равен единице, поэтому точка не входит. Граничные точки дают ровно одну вторую и поэтому включаются.

\textbf{\result}
\[
\boxed{x\in\left[\frac\pi3+2\pi n;\frac{5\pi}3+2\pi n\right],\quad n\in\mathbb Z.}
\]

\textbf{Задача на закрепление}

Решите неравенство
\[
\cos x\le -\frac12.
\]
\end{document}
